% =====================================================================
%  P2.  No weight extracts C(N): a no-go for the divisor-switch route,
%       with the exact identities of its design space.
%
%  Deployment version.  Written to deploy/SPEC.md.
%  Compile:  pdflatex P2-no-go-divisor-switch.tex   (twice)
% =====================================================================
\documentclass[11pt,a4paper]{article}

\usepackage{amsmath}
\usepackage{amssymb}
\usepackage{amsthm}
\usepackage[margin=1in]{geometry}
\usepackage[colorlinks=true,linkcolor=blue,citecolor=blue]{hyperref}

\theoremstyle{plain}
\newtheorem{theorem}{Theorem}
\newtheorem{corollary}[theorem]{Corollary}
\newtheorem{proposition}[theorem]{Proposition}
\newtheorem{lemma}[theorem]{Lemma}
\theoremstyle{definition}
\newtheorem{observation}[theorem]{Observation}
\newtheorem{measurement}[theorem]{Measurement}
\theoremstyle{remark}
\newtheorem{note}[theorem]{Note}

\renewcommand{\SS}{\mathfrak{S}}
\newcommand{\AAA}{\mathfrak{A}}
\newcommand{\Emu}{E_\mu}
\DeclareMathOperator{\rad}{rad}
\DeclareMathOperator{\lcm}{lcm}

\title{No weight extracts $\sum_{n<N}\Lambda(n)\mu(N-n)$:\\
a no-go for the divisor-switch route,\\
with the exact identities of its design space}
\author{}
\date{}

\begin{document}
\maketitle

\begin{abstract}
Let $N$ be a large even integer and
$C(N)=\sum_{n<N}\Lambda(n)\mu(N-n)$. The Huang--Li reduction of binary
Goldbach to Elliott--Halberstam \cite{HL} consumes its hypothesis
through a M\"obius-weighted correlation sum in the fixed class
$n\equiv N\ (k)$, carrying a weight $w_k$; the two weights that occur
are $w_k=1$, for which the sum is unconditionally small, and
$w_k=\log k$, for which it is equivalent to Goldbach itself
\cite{P1}. It is natural to ask whether some weight sits between them:
one whose complete divisor sum is still cheap but whose main-term
coefficient is still of size $1$, and which would therefore extract
$C(N)$ from Bombieri--Vinogradov alone.

We show that no such weight exists. Writing $b=\mu*w$, the two
requirements pull in opposite directions --- extraction needs mass in
$b$ at the truncation point $K=N^{\theta'}$
(Lemma~\ref{lem:extract}), while the residual is accessible to
Bombieri--Vinogradov only for $b$ supported below $N^{\theta'-1/2}$
(Lemma~\ref{lem:bv}) --- and the two thresholds are separated by a
factor $N^{1/2}$. Theorem~\ref{thm:D} converts this into the bound
$\|b\|_1/|B_w|\gg\exp(c_1\sqrt{(1/2)\log N})$, which exceeds every
power of $\log N$; Theorem~\ref{thm:Dprime} shows the separation
survives granting $EH(N^{\theta_E})$ for any fixed $\theta_E<1$, so
the obstruction is not a shortage of level. The most natural family of
weights, $w_k=f(\log k)$ with $f$ polynomial, falls outside the
support hypothesis and is closed separately
(Proposition~\ref{prop:Dpp}): its complete part is a nonnegative
Chen-type sum, never $o(N)$, and the tuning that would cancel it is
circular. For completeness we record the exact position of the other
classical mechanism: the circle method applied directly to $C(N)$ has
zero margin (Proposition~\ref{prop:E}), both standard pairings lying
at or above the trivial bound, with Parseval supplying the floor.

The design space is first laid out exactly. Seven finite rearrangements
(Propositions~\ref{prop:dilate}--\ref{prop:direct}) identify the
per-modulus discrepancy as a dilate of a M\"obius--prime correlation,
show that the signs $\mu(k)$ cancel out of the $\log k$ branch
exactly, and exhibit the Goldbach count as a signed sum of
nonnegative layers. These identities are what make the no-go a
statement about a space rather than about two examples.

Net progress toward Goldbach: zero. The value of a no-go of this kind
is that it closes a design space rather than leaving it to be
re-explored.
\end{abstract}

% =====================================================================
\section{Introduction}

\subsection{The setting}

Let $N$ be a large even integer, $\theta'\in(1/2,1)$ fixed, and
$K=N^{\theta'}$. Put
\begin{equation}\label{eq:C}
  C(N)=\sum_{n<N}\Lambda(n)\mu(N-n),
  \qquad
  A(N;k)=\sum_{\substack{n<N\\ n\equiv N\ (k)}}\Lambda(n)\mu(N-n),
\end{equation}
and for a weight $w:\mathbb N\to\mathbb R$,
\begin{equation}\label{eq:Tw}
  T_w \;=\;\sum_{\substack{k<K\\ (k,N)=1}}\mu(k)\,w_k
     \Bigl[A(N;k)-\frac{C(N)}{\varphi(k)}\Bigr].
\end{equation}
The functionals $E_4$ and $E_3$ of \cite{HL} are $T_w$ with $w_k=1$
and $w_k=\log k$ respectively. The first is unconditionally
$\ll_A N(\log N)^{-A}$; the second is, by an unconditional identity,
equivalent to Huang--Li's equation~(22) and hence already gives binary
Goldbach \cite{P1}. This paper closes the space between them.

Throughout, $\varphi$ is Euler's function, $\rad(u)=\prod_{p\mid u}p$,
$c,c_1,\dots$ are positive absolute constants, and
\begin{equation}\label{eq:AN}
  \AAA(N)=\prod_{q\,\nmid\,N}\Bigl(1-\frac{1}{q(q-1)}\Bigr),
  \qquad
  \SS(N)=2\prod_{p>2}\Bigl(1-\tfrac{1}{(p-1)^2}\Bigr)
          \prod_{\substack{p\mid N\\ p>2}}\Bigl(1+\tfrac{1}{p-2}\Bigr).
\end{equation}

\subsection{The design space}

Let $w$ be arbitrary and let $b:=\mu*w$, so that
$w_k=\sum_{d\mid k}b_d$; every weight has this form, and $b$ is
determined by $w$. Put
\begin{equation}\label{eq:Bw}
  B_w:=\sum_{\substack{k<K\\(k,N)=1}}\frac{\mu(k)w_k}{\varphi(k)},
  \qquad
  \|b\|_1:=\sum_{u<N}|b_u| .
\end{equation}
Running the divisor switch of \cite{P1} with the weight $w$ gives the
identity
\begin{equation}\label{eq:extract}
  B_w\cdot C(N)\;=\;
  \underbrace{\sum_{u<N}\Lambda(N-u)\mu^2(u)\,b_u}_{\text{complete part}}
  \;-\;\underbrace{\mathcal R_w}_{\text{residual}}
  \;+\;O\!\left(N^{o(1)}\right),
\end{equation}
the complete part being evaluated by Lemma~\ref{lem:Gb} below, which
identifies $b$ as exactly the complete divisor transform.

\begin{lemma}\label{lem:Gb}
For squarefree $u$, $\sum_{k\mid u}\mu(k)w_k=\mu(u)\,b_u$.
\end{lemma}

\begin{proof}
$\sum_{k\mid u}\mu(k)\sum_{d\mid k}b_d
 =\sum_{d\mid u}b_d\,\mu(d)\sum_{j\mid(u/d)}\mu(j)
 =\sum_{d\mid u}b_d\,\mu(d)\,[\,u/d=1\,]=\mu(u)b_u$,
using $(d,u/d)=1$ for squarefree $u$.
\end{proof}

The two known cases are the two ends of this space: $w=1$ gives
$b=\delta_1$, $\|b\|_1=1$ and (Huang--Li's Lemma~1)
$B_w\ll e^{-c\sqrt{\log K}}$; while $w=\log$ gives $b=\Lambda$,
$\|b\|_1\asymp N$ --- the complete part \emph{is} the binary Goldbach
sum --- and $B_w\to-\SS(N)\asymp1$.

\subsection{Results}

Section~\ref{sec:identities} records the exact structure of the space:
the per-modulus discrepancy is a dilate (Proposition~\ref{prop:dilate}),
the signs $\mu(k)$ cancel exactly on the $\log k$ branch
(Proposition~\ref{prop:posweights}), both known weights are the same
sum of dilated walls with different weights
(Proposition~\ref{prop:flatsum}), the untruncated object is the
Goldbach count itself (Proposition~\ref{prop:untrunc}), which is
therefore a signed sum of nonnegative layers
(Proposition~\ref{prop:layers}) over an ordinary
Bombieri--Vinogradov modulus (Proposition~\ref{prop:combined}); and
the wall cancels out of the count entirely
(Proposition~\ref{prop:direct}).

Section~\ref{sec:nogo} proves the no-go: Lemmas~\ref{lem:extract}
and~\ref{lem:bv} are the two opposing constraints,
Theorem~\ref{thm:D} is their consequence, and
Theorem~\ref{thm:Dprime} shows it survives Elliott--Halberstam.
Section~\ref{sec:Dpp} closes polynomial weights, which lie outside
the hypothesis of Theorem~\ref{thm:D}. Section~\ref{sec:circle}
records the position of the circle method.

\subsection{What is not claimed}\label{sec:notclaimed}

\begin{itemize}
\item \textbf{This is a no-go for one method.} Theorem~\ref{thm:D}
closes divisor switching with Bombieri--Vinogradov as the only input,
over the whole weight space of that method. It is not, and does not
claim to be, an obstruction to other methods --- in particular it says
nothing about the spectral large sieve of \cite{Li23}.

\item \textbf{The genre is classical.} Bombieri's asymptotic sieve
\cite{Bom76} shows that sieve weights alone cannot detect primes,
however chosen, and the parity obstruction it isolates is the reason
one expects a statement of this shape. What is added here is the
precise quantitative form for this route
(Note~\ref{note:bombieri}).

\item \textbf{No progress toward Goldbach.} Nothing here bounds
$C(N)$, and nothing here is a step toward bounding it.

\item \textbf{The identities of Section~\ref{sec:identities} are
rearrangements.} They are exact and they are finite; none of them is
an estimate, and none of them makes any quantity smaller.
\end{itemize}

% =====================================================================
\section{The design space, exactly}\label{sec:identities}

\begin{proposition}[the discrepancy is a dilate]\label{prop:dilate}
Unconditionally,
\begin{equation}\label{eq:dilate}
  A(N;k)\;=\;\mu(k)\,H(N;k),
  \qquad
  H(N;k)\;=\;\sum_{\substack{m<N/k\\ (m,k)=1}}\Lambda(N-mk)\,\mu(m).
\end{equation}
\end{proposition}

\begin{proof}
The class $n\equiv N\pmod k$ is $k\mid N-n$, so writing $N-n=mk$ turns
the sum into $\sum_{m<N/k}\Lambda(N-mk)\mu(mk)$; and
$\mu(mk)=\mu(m)\mu(k)$ when $(m,k)=1$, zero otherwise.
\end{proof}

Identity~\eqref{eq:dilate} qualifies the shape of the whole problem.
$EH_\mu$ is consumed where one M\"obius argument \emph{divides} the
other, but the per-modulus discrepancy is $\mu(k)$ times a sum in
which the two arguments are related by a \emph{difference}, at scale
$N/k$. The two couplings are not two fields; they are the same field
seen before and after the switch. It also explains why $\mu$ is the
noisier of the two inputs in progressions: for $\mu$,
$|A(N;k)|=|H(N;k)|$ is a M\"obius--prime correlation of length $N/k$
and nothing makes it small, whereas for a random sign pattern
$\varepsilon$ the sum $\sum_m\Lambda(N-mk)\varepsilon(mk)$ is a sum of
independent signs and exhibits square-root cancellation.

\begin{proposition}[the weights are nonnegative]\label{prop:posweights}
Substituting \eqref{eq:dilate} into $T_{\log}$ and using $\mu(k)^2=1$
on squarefree $k$,
\begin{equation}\label{eq:posweights}
  T_{\log}\;=\;\sum_{\substack{k<K\\(k,N)=1}}
    \mu^2(k)\,(\log k)\,H(N;k)\;-\;C(N)\,B_{\log}(K),
  \qquad
  B_{\log}(K)=\!\!\sum_{\substack{k<K\\(k,N)=1}}\!\!
     \frac{\mu(k)\log k}{\varphi(k)} .
\end{equation}
\end{proposition}

\begin{proof}
Immediate from \eqref{eq:dilate} and $\mu(k)\mu(k)=\mu^2(k)$.
\end{proof}

The signs $\mu(k)$ cancel exactly: they cancel against the $\mu(k)$
that \eqref{eq:dilate} extracts from the progression sum, and what is
left is a sum of dilated M\"obius--prime correlations with
\emph{nonnegative} weights $\mu^2(k)\log k$. Huang--Li discard the
$\mu(k)$ by the triangle inequality; on the $w_k=1$ branch keeping
them is what makes the unconditional theorem of \cite{P1} possible; on
the $\log k$ branch they are not there to keep. Whatever smallness
$T_{\log}$ has must therefore come from the signs of $H(N;k)$ across
$k$.

\begin{proposition}[both ends are the same sum of dilated walls]%
\label{prop:flatsum}
For any weight $w$ that is a function of $k$ alone,
\begin{equation}\label{eq:flatsum}
  T_w\;=\;\sum_{\substack{k<K\\(k,N)=1}}w_k\,H(N;k)\;-\;B_w\,C(N).
\end{equation}
In particular
\[
  T_1=\!\!\sum_{\substack{k<K\\(k,N)=1}}\!\!H(N;k)-B_1C(N),
  \qquad
  T_{\log}=\!\!\sum_{\substack{k<K\\(k,N)=1}}\!\!(\log k)H(N;k)-B_{\log}C(N),
\]
and both weights are nonnegative on the range.
\end{proposition}

\begin{proof}
Put \eqref{eq:dilate} into \eqref{eq:Tw}: for squarefree $k$,
$\mu(k)A(N;k)=\mu(k)^2H(N;k)=H(N;k)$, and the subtracted mean term
assembles into $B_wC(N)$ by \eqref{eq:Bw}.
\end{proof}

Since $B_1\ll e^{-c\sqrt{\log K}}$ kills the second term while
$B_{\log}\to-\SS(N)$ does not, the unconditional theorem of \cite{P1}
says exactly that the \emph{flat} sum of dilated walls is
$\ll_AN(\log N)^{-A}$, and the Goldbach problem says the same sum
weighted by $\log k$ is not small. Everything between what is proved
and what is open is the factor $\log k$ inside a positively weighted
sum of the same terms.

\begin{measurement}[the two ends, measured]\label{meas:flatsum}
Over $N=2\cdot10^5$ to $2.56\cdot10^7$ by doubling, at
$\theta'=0.56$, \eqref{eq:flatsum} holds to a worst relative error of
$1.875\cdot10^{-16}$, and
\[
  \frac{\bigl|\sum_kH(N;k)\bigr|}{\bigl|\sum_k(\log k)H(N;k)\bigr|}
  \;=\;0.1807,\ 0.1740,\ 0.1624,\ 0.1389,\ 0.1456,\ 0.1216,\
       0.1258,\ 0.1188,
\]
falling throughout; at the top $|T_1|/N=0.01425$ against
$|T_{\log}|/N=0.1245$. Fitted against $\log N$ the flat sum decays as
$N^{-0.3620}$ and $|T_{\log}|/N$ as $N^{-0.2658}$: the gap widens, as
it must if one side is $\ll_AN(\log N)^{-A}$ and the other is
$\asymp N$. Identity~\eqref{eq:dilate} itself is verified over every
squarefree $k<N^{0.56}$ coprime to $N$ at
$N=2\cdot10^5,\,4\cdot10^5,\,8\cdot10^5$ with worst relative error
$2.294\cdot10^{-14}$, $1.197\cdot10^{-13}$, $4.591\cdot10^{-14}$, and
\eqref{eq:posweights} against a direct evaluation of $T_{\log}$ to at
most $1.8\cdot10^{-16}$ over $N=2\cdot10^5$ to $3.2\cdot10^6$
(\texttt{lab\_weight\_gap.py}, \texttt{lab\_dilate\_identity.py},
\texttt{lab\_positive\_weights.py}).
\end{measurement}

\begin{proposition}[the untruncated sum is the count]%
\label{prop:untrunc}
Summing over \emph{all} $k$, with no truncation and no coprimality
restriction,
\begin{equation}\label{eq:untrunc}
  \sum_{k\ge1}(\log k)\,\mu(k)\,A(N;k)
   \;=\;-\sum_{n<N}\Lambda(n)\,\mu(N-n)\Lambda(N-n)
   \;=\;\sum_{p<N}\Lambda(N-p)\log p .
\end{equation}
So the untruncated object \emph{is} $\tilde r(N)$, up to its
prime-power part.
\end{proposition}

\begin{proof}
Both exchanges are finite sums. Since $A(N;k)$ sums $n$ over
$k\mid N-n$, exchanging the order and applying
$\sum_{d\mid u}\mu(d)\log d=-\Lambda(u)$ --- M\"obius inversion of
$\log=\mathbf 1*\Lambda$ --- gives the first equality. The second
holds because $\mu(u)\Lambda(u)=-\log p$ at $u=p$ prime and vanishes
at every higher prime power.
\end{proof}

Proposition~\ref{prop:untrunc} is not a decoration: it says that
\eqref{eq:flatsum} with $w=\log$ is a \emph{rearrangement} of
$\tilde r(N)$ rather than an approximation waiting for an estimate,
so all of its content is in the truncation at $K$.

\begin{proposition}[the count as signed layers]\label{prop:layers}
Cutting the same double sum along the inner variable of
\eqref{eq:dilate},
\begin{equation}\label{eq:layers}
  \sum_{k}(\log k)H(N;k)\;=\;\sum_{m}\mu(m)\,L(N;m),
  \qquad
  L(N;m):=\!\!\sum_{\substack{k<N/m,\ (k,m)=1\\ \mu(k)\neq0}}\!\!
            (\log k)\,\Lambda(N-mk),
\end{equation}
and every $L(N;m)$ is nonnegative. With
Proposition~\ref{prop:untrunc}, the Goldbach count is therefore an
alternating sum, signed by $\mu$, of nonnegative layers.
\end{proposition}

\begin{proof}
A finite exchange; the coprimality conditions on the pair $(k,m)$ are
symmetric. Nonnegativity is that of $\log k$ on $k\ge1$ and of
$\Lambda$.
\end{proof}

\begin{proposition}[the count over a combined modulus]%
\label{prop:combined}
Expanding the squarefree condition inside \eqref{eq:layers} by
$\mu^2(k)=\sum_{d^2\mid k}\mu(d)$ and writing $k=d^2j$,
\begin{equation}\label{eq:combined}
  \sum_{p<N}\Lambda(N-p)\log p
   \;=\!\!\sum_{\substack{m,d\ \text{squarefree}\\ (d,m)=1,\ md^2<N}}\!\!
     \mu(m)\mu(d)\,L_2(N;m,d),
  \qquad
  L_2(N;m,d)=\!\!\sum_{\substack{j<N/(md^2)\\ (j,m)=1}}\!\!
   \log(d^2j)\,\Lambda\bigl(N-md^2j\bigr),
\end{equation}
and $L_2$ carries no squarefree condition: it is a prime count in the
progression $p\equiv N\pmod{md^2}$, an unadorned
Bombieri--Vinogradov object of modulus $q=md^2$.
\end{proposition}

\begin{proof}
Finite exchange, and $(k,m)=1$ splits as $(d,m)=(j,m)=1$.
\end{proof}

\begin{proposition}[the wall cancels out of the count]%
\label{prop:direct}
Substituting \eqref{eq:posweights} into the identity of \cite[Thm.~C]{P1},
the two occurrences of $C(N)$ carry opposite signs and cancel:
\begin{equation}\label{eq:direct}
  \tilde r(N)\;=\;\SS(N)N
   \;+\!\!\sum_{\substack{k<K\\(k,N)=1}}\!\!(\log k)H(N;k)
   \;-\;C(N)\bigl(B_{\log}(K)+\SS(N)\bigr)
   \;+\;O_A\!\left(\frac{N}{(\log N)^{A}}\right).
\end{equation}
Since $B_{\log}(K)\to-\SS(N)$, the quantity $C(N)$ enters the Goldbach
count only through the vanishing factor $B_{\log}(K)+\SS(N)$.
Consequently
\begin{equation}\label{eq:directcond}
  \sum_{\substack{k<K\\(k,N)=1}}(\log k)\bigl|H(N;k)\bigr|
   \;\le\;(1-\varepsilon)\,\SS(N)\,N
\end{equation}
suffices for $\tilde r(N)>0$, for all large even $N$.
\end{proposition}

\begin{proof}
Immediate from \eqref{eq:direct} and $|C(N)|\le\AAA(N)N(1+o(1))$, the
last bound now multiplied by $B_{\log}(K)+\SS(N)=o(1)$ rather than by
$\SS(N)$.
\end{proof}

\begin{note}[what the cancellation buys, and what it does not]
Condition~\eqref{eq:directcond} has the same shape as the
constant-factor condition of \cite[Prop.~5]{P1}: a bound at level
$N^{\theta'}$, with no saving in $\log N$. Nothing has moved off the
level axis. What has changed is the threshold constant, from
$\SS(N)(1-\AAA(N))$ to $\SS(N)$, and the change removes the arithmetic
sensitivity along with the slack, because the two were the same thing:
$1-\AAA(N)$ collapses at primorials, and $\SS(N)$ does not. This is
offered as a weaker sufficient condition, proved from two identities.
It is not offered as something computation confirms: at accessible $N$
the unspecified $O_A$ term of \eqref{eq:direct} is numerically the
very quantity \eqref{eq:directcond} must bound.
\end{note}

\begin{measurement}[the two thresholds, compared]\label{meas:direct}
At $\theta'=0.56$ over $N=2\cdot10^5$ to $3.2\cdot10^6$, with
$B_H(N)=\sum_{k<K}(\log k)|H(N;k)|$, the new ratio
$B_H(N)/(\SS(N)N)$ reads
$0.4578,\,0.4064,\,0.4079,\,0.3769,\,0.3338$ --- under $1$ throughout
and falling --- against
$2.1591,\,1.9747,\,1.9500,\,1.7483,\,1.5798$ for the old ratio
$B_H(N)/(\SS(N)(1-\AAA(N))N)$, which is above $1$ throughout. Across a
seven-point family of arithmetically diverse $N$ the new ratio spans a
factor $2.00$ where the old spans $23.25$. The residual coupling is
negligible at these $N$: $|B_{\log}+\SS|/\SS$ stays under $0.0383$ and
$|C(N)(B_{\log}+\SS)|/N$ never exceeds $3.315\cdot10^{-4}$, four
orders below $\SS(N)$. The residual of \eqref{eq:direct} decays as
$N^{-0.2794}$ with correlation $0.99930$ and reaches $2\%$ of $\SS$
only near $N=10^{10.16}$
(\texttt{lab\_direct\_route.py}, \texttt{lab\_direct\_identity.py}).
\end{measurement}

% =====================================================================
\section{The no-go}\label{sec:nogo}

\subsection{The two opposing constraints}

\begin{lemma}[extraction]\label{lem:extract}
$\displaystyle B_w=\sum_{\substack{d<K\\(d,N)=1}}b_d\,
 \frac{\mu(d)}{\varphi(d)}\,\rho_{dN}\!\left(\frac{K}{d}\right)$,
where $\rho_{n}(x)=\sum_{j<x,\,(j,n)=1}\mu(j)/\varphi(j)
 \ll e^{-c_1\sqrt{\log x}}$.
\end{lemma}

\begin{proof}
Write $k=dj$; for squarefree $k$ one has $(d,j)=1$,
$\mu(k)=\mu(d)\mu(j)$ and $\varphi(k)=\varphi(d)\varphi(j)$. The bound
on $\rho$ is Huang--Li's Lemma~1, a form of
Goldston--Y{\i}ld{\i}r{\i}m \cite{GY}.
\end{proof}

\begin{lemma}[BV-accessibility]\label{lem:bv}
Expanding $w_k=\sum_{d\mid k}b_d$ inside the residual gives
$\mathcal R_w=\sum_db_d\,\mathcal R^{(d)}$, where $\mathcal R^{(d)}$
is a sum of $\Lambda$ over progressions to moduli $md$ with
$m<N^{1-\theta'}$. Bombieri--Vinogradov bounds each
$\mathcal R^{(d)}\ll_AN(\log N)^{-A}$ provided
$d\le N^{1/2-\delta}/N^{1-\theta'}=N^{\theta'-1/2-\delta}$; hence, on
that support, $\mathcal R_w\ll_A\|b\|_1\,N(\log N)^{-A}$.
\end{lemma}

\begin{proof}
The residual of the switch is \cite[eq.~(R1)]{P1}: after unfolding
$\mu^2$ and the coprimality condition, every term is a sum of
$\Lambda$ over a progression to modulus $q=m\lcm(d^2,e)$ with
$m<N^{1-\theta'}$ and $d,e$ at most a fixed power of $\log N$;
Bombieri--Vinogradov is applied in the form of \cite{MV}.
Inserting the extra divisor $d\mid k$ multiplies the modulus by $d$,
and the level condition $q\le N^{1/2-\delta}$ becomes the stated
bound on $d$. Summing the $\ll_AN(\log N)^{-A}$ over the support of
$b$ with the trivial weight $|b_d|$ gives the claim.
\end{proof}

The two lemmas pull in opposite directions. Lemma~\ref{lem:extract}
says $B_w$ is controlled by the part of $b$ sitting \emph{at} the
truncation point $K$: mass at $d\le K^{1-\varepsilon}$ is damped by
$e^{-c\sqrt{\varepsilon\log K}}$. Lemma~\ref{lem:bv} says $b$ may only
live \emph{below} $N^{\theta'-1/2}$. The two thresholds are separated
by a factor $N^{1/2}$.

\subsection{The theorem}

\begin{theorem}[no weight extracts $C(N)$]\label{thm:D}
Fix $\theta'\in(1/2,1)$ and $\delta>0$, and let $w$ be any weight
whose transform $b=\mu*w$ is supported in
$[1,\,N^{\theta'-1/2-\delta}]$ --- the range in which the residual is
accessible to Bombieri--Vinogradov. Then
\[
  \frac{\|b\|_1}{|B_w|}\;\gg\;
   \exp\!\left(c_1\sqrt{(\tfrac12+\delta)\log N}\right),
\]
and consequently \eqref{eq:extract} yields at best
\[
  |C(N)|\;\ll_A\;\exp\!\left(c_1\sqrt{\tfrac12\log N}\right)
   \cdot\frac{N}{(\log N)^{A}},
\]
which is not a saving of any power of $\log N$. No weight in the
design space extracts $C(N)$ by divisor switching plus
Bombieri--Vinogradov.
\end{theorem}

\begin{proof}
By Lemma~\ref{lem:extract} and $\varphi(d)\ge1$,
\[
  |B_w|\le\sum_{d\le D}\frac{|b_d|}{\varphi(d)}
    \left|\rho_{dN}\!\left(\frac Kd\right)\right|
  \le\|b\|_1\max_{d\le D}\left|\rho_{dN}\!\left(\frac Kd\right)\right|
  \ll\|b\|_1\,e^{-c_1\sqrt{\log(K/D)}},
\]
with $D=N^{\theta'-1/2-\delta}$, so that $K/D=N^{1/2+\delta}$. This is
the displayed ratio. Feeding it into \eqref{eq:extract} together with
the complete-part bound $\ll\|b\|_1\log N$ (Lemma~\ref{lem:Gb}) and
Lemma~\ref{lem:bv} gives the bound on $|C(N)|$; and
$e^{c_1\sqrt{\log N/2}}$ exceeds every power of $\log N$.
\end{proof}

The loss factor is $\exp(c_1\sqrt{(1/2)\log N})$, and the $1/2$ in the
exponent is literally the exponent of the $\sqrt N$ barrier: it is the
gap between where a weight must live to see $C(N)$ (at $N^{\theta'}$,
the truncation point) and where it may live for
Bombieri--Vinogradov to close its residual (below $N^{\theta'-1/2}$).
The two known weights exhibit the two endpoints of this space;
Theorem~\ref{thm:D} shows the interior is empty, and shows why.

\begin{note}[relation to Bombieri's asymptotic sieve]\label{note:bombieri}
The genre of Theorem~\ref{thm:D} is classical. Bombieri's asymptotic
sieve \cite{Bom76} shows that sieve weights alone cannot detect
primes, however the weights are chosen, and the parity obstruction it
isolates is the reason one expects a statement of this shape. What
Theorem~\ref{thm:D} adds is not the phenomenon but its precise form
here: the obstruction is quantified as a separation of $N^{1/2}$
between two thresholds, it is specific to the divisor-switch route,
and it survives granting the full Elliott--Halberstam conjecture
rather than merely Bombieri--Vinogradov. The reason the statement is
worth making is narrow but real: without it, the $\log k$ weight looks
like one member of a family of candidate consumptions, and the
identity of \cite[Thm.~C]{P1} then looks like an accident of one
choice of weight rather than a property of the route.
\end{note}

\subsection{The no-go survives Elliott--Halberstam}

The obvious objection to Theorem~\ref{thm:D} is that
Bombieri--Vinogradov is not the only available input: Huang--Li's
Theorem~1 itself \emph{assumes} $EH$ for $\Lambda$ at level
$N^{\theta}$. Raising the assumed level does not help.

\begin{theorem}\label{thm:Dprime}
Suppose $\Lambda$ has level of distribution $\theta_E\in(0,1)$, i.e.
$EH(N^{\theta_E})$ holds. Then the residual is accessible only for $b$
supported on $d\le N^{\theta_E-(1-\theta')}$, while extraction still
requires mass at $d\asymp K=N^{\theta'}$; the two thresholds are
separated by $N^{1-\theta_E}$, and
\[
  \frac{\|b\|_1}{|B_w|}\;\gg\;
  \exp\!\left(c_1\sqrt{(1-\theta_E)\log N}\right),
\]
which exceeds every power of $\log N$ for each fixed $\theta_E<1$.
\end{theorem}

\begin{proof}
Identical to Theorem~\ref{thm:D} with $N^{1/2-\delta}$ replaced by
$N^{\theta_E}$: the residual's moduli are $md$ with $m<N^{1-\theta'}$,
so $d\le N^{\theta_E}/N^{1-\theta'}$, whence
$K/D=N^{\theta'}/N^{\theta_E-1+\theta'}=N^{1-\theta_E}$, and
Lemma~\ref{lem:extract} gives the ratio.
\end{proof}

Closing the gap would require $\theta_E=1$ exactly: equidistribution
of $\Lambda$ in progressions to moduli of size $N$ itself, where each
progression contains $O(1)$ terms and the statement carries no
information. So the route stays closed \emph{even granting the full
Elliott--Halberstam conjecture} --- it needs not a stronger level but
a different mechanism.

\begin{measurement}[the load-bearing lemma, checked as an identity]%
\label{meas:extract}
Lemma~\ref{lem:extract} carries the theorem, so it is checked
directly at $N=99\,999\,998$, $\theta'=0.56$,
$K=\lfloor N^{\theta'}\rfloor=30199$. For the single-divisor weights
$w_k=[\,d_0\mid k\,]$, i.e. $b=\delta_{d_0}$, the lemma is an identity
and is verified as one: over all $10\,262$ squarefree $d_0<K$ coprime
to $N$, the brute-force $B_w$ over $k<K$ and the factorised
$\mu(d_0)\varphi(d_0)^{-1}\rho_{d_0N}(K/d_0)$ agree to
$6.5\cdot10^{-18}$.

The size behaves as the lemma predicts. The quantity
$|B_w|\varphi(d_0)$ is of order $1$ only when $K/d_0=O(1)$, i.e. only
when the weight's transform sits at the truncation point, and is
damped throughout the range Bombieri--Vinogradov admits: for these
parameters $N^{\theta'-1/2}=3$, and over $d_0\le3$ the largest value
attained is $0.004145$. The value is not a function of $K/d_0$ alone
--- $\rho_{d_0N}$ carries the condition $(j,d_0N)=1$, so it depends on
which small primes divide $d_0$. At $\lfloor K/d_0\rfloor=7$ the $184$
admissible $d_0$ produce exactly four values, split by $\gcd(d_0,15)$:
\[
\begin{array}{r|cccc}
 \gcd(d_0,15) & 1 & 3 & 5 & 15\\\hline
 |B_w|\,\varphi(d_0) & 0.250000 & 0.750000 & 0.500000 & 1.000000
\end{array}
\]
and at $\lfloor K/d_0\rfloor=13$ the $55$ admissible $d_0$ produce six
values spanning $0.066667$ to $0.566667$
(\texttt{audit\_extraction\_tradeoff.py}).
\end{measurement}

% =====================================================================
\section{Smooth weights: $\mu*\log^D=\Lambda_D$}\label{sec:Dpp}

Theorem~\ref{thm:D} closes the weights whose transform $b=\mu*w$ is
supported low enough for Bombieri--Vinogradov. That hypothesis
excludes the most natural family of all, $w_k=f(\log k)$ with $f$ a
polynomial, whose transform is spread over every scale. For those the
structure is explicit:
\[
  b\;=\;\mu*\log^D\;=\;\Lambda_D,
\]
the generalised von Mangoldt function. $\Lambda_D$ vanishes on
integers with more than $D$ prime factors, and on a squarefree
$u=p_1\cdots p_r$ with $r\le D$ it is the $r$-th finite difference of
$x^D$ at the points $\log p_i$; in particular
$\Lambda_D(u)=D!\,\log p_1\cdots\log p_D$ when $r=D$. Hence for
$f=x^D$ the complete part splits by the number of prime factors,
\begin{equation}\label{eq:CPD}
  \mathrm{CP}_D(N)=\sum_{r=1}^{D}\
    \sum_{\substack{u<N\ \text{squarefree}\\ \omega(u)=r}}
   \Lambda(N-u)\,\Lambda_D(u),
\end{equation}
the $r=1$ piece being a Goldbach-type sum and the $r\ge2$ pieces
Chen-type sums (prime plus $r$-almost-prime).

\begin{proposition}[polynomial weights]\label{prop:Dpp}
Let $w_k=f(\log k)$ with $f=\sum_{j\le D}c_jx^j$ real.
\begin{enumerate}
\item[(i)] If $D=0$ then $B_w\ll e^{-c\sqrt{\log K}}$ (Huang--Li's
Lemma~1): no extraction.
\item[(ii)] If $f$ is a single monomial $x^D$, $D\ge1$, then every
term of $\mathrm{CP}_D$ is nonnegative, since $\Lambda\ge0$ and
$\Lambda_D\ge0$; by classical sieve bounds
$\mathrm{CP}_D\asymp N(\log N)^{D-1}$, with a fixed sign. It is never
$o(N)$.
\item[(iii)] Cancellation therefore requires coefficients of opposite
sign, i.e. tuning $c_1$ against $c_2,\dots$; but the ratio that would
have to be matched involves the asymptotics of the $r=1$ piece, which
for $D=1$ \emph{is} the binary Goldbach sum. The tuning is circular.
\end{enumerate}
Consequently no polynomial weight makes the complete part a
classically known object.
\end{proposition}

\begin{proof}
(i) is Huang--Li's Lemma~1 applied to $w\equiv1$. For (ii),
$\Lambda_D\ge0$ on the squarefree integers with $\omega\le D$ because
each $\Lambda_D(u)$ is a finite difference of the convex increasing
function $x\mapsto x^D$ at $\log p_1<\cdots<\log p_r$; every summand of
\eqref{eq:CPD} is then a product of nonnegative factors, so no
cancellation is available and the order of magnitude is that of the
sieve upper and lower bounds for prime-plus-almost-prime sums. For
(iii), a degree-$D$ tuning has $D$ free coefficients and the
constraint that must be solved is an identity between the leading
asymptotics of the $r$-pieces; the $r=1$ piece with $D=1$ is
$\sum_{p<N}\Lambda(N-p)\log p$, whose asymptotic is the assertion to
be proved.
\end{proof}

\begin{note}[the canonical tuning moves the wrong way]
Since $b=\mu*w$ we have $w=\mathbf 1*b$ and hence
$B(s)=W(s)/\zeta(s)$. That series has a \emph{double} pole at $s=1$
for every degree-$2$ $f$, so no such $f$ removes it and $b$ cannot be
made mean-zero; the most a degree-$2$ tuning can do is kill the
second-order term. With $f=x^2+cx$ one has $W=\zeta''-c\zeta'$, and
neither $\zeta''$ nor $\zeta'$ carries a $(s-1)^{-1}$ term.
\end{note}

\begin{measurement}[the two pieces are the same order]\label{meas:poly}
Nonnegativity, and not a separation of scales, is what closes (ii).
At $N=10^6,\,4\cdot10^6,\,1.6\cdot10^7$ the two pieces of
$\mathrm{CP}_2$ are the same order:
\[
\begin{array}{r|ccc}
 N & 10^6 & 4\cdot10^6 & 1.6\cdot10^7\\\hline
 r=1:\ \sum_p\Lambda(N-p)\log^2p\,/\,N & 22.5145 & 25.0438 & 27.4579\\
 r=2:\ 2\sum_{pq}\Lambda(N-pq)\log p\log q\,/\,N
   & 17.3648 & 19.7926 & 22.2512\\
 \text{ratio } r_2/r_1 & 0.7713 & 0.7903 & 0.8104\\
 \mathrm{CP}_2/(N\log N) & 2.8866 & 2.9494 & 2.9967
\end{array}
\]
The ratio drifts toward $1$, not toward $0$ or $\infty$. Calibration:
the $D=1$ column reproduces the Goldbach sum
$\sum_p\Lambda(N-p)\log p=1.7565N,\,1.7633N,\,1.7614N$ against
$\SS(N)=1.76043$, which is the same at all three $N$ because
$N$ has the same prime support $\{2,5\}$ in each case
(\texttt{audit\_polyweight.py}).
\end{measurement}

% =====================================================================
\section{The circle method has zero margin on $C(N)$}\label{sec:circle}

Since the switch is closed over its whole design space, it is worth
recording the exact position of the other classical mechanism. Write
\[
  C(N)=\int_0^1S_\Lambda(\alpha)\,S_\mu(-\alpha)\,e(-N\alpha)\,d\alpha,
  \qquad
  S_\Lambda(\alpha)=\sum_{n\le N}\Lambda(n)e(n\alpha),\quad
  S_\mu(\alpha)=\sum_{m\le N}\mu(m)e(m\alpha).
\]

\begin{proposition}[zero margin]\label{prop:E}
Both standard ways of estimating this integral lie at or above the
trivial bound $|C(N)|\le\psi(N)\sim N$:
\begin{enumerate}
\item[(i)] $\|S_\Lambda\|_2\|S_\mu\|_2
   =\bigl(\sum_{n\le N}\Lambda(n)^2\bigr)^{1/2}
     \bigl(\sum_{m\le N}\mu^2(m)\bigr)^{1/2}
   \sim(6/\pi^2)^{1/2}\,N(\log N)^{1/2}$, exceeding the trivial bound
   by a factor $\asymp(\log N)^{1/2}$, which \emph{grows};
\item[(ii)] any bound of the shape
  $\sup_\alpha|S_\mu(\alpha)|\cdot\|S_\Lambda\|_1$ is at least
  $\|S_\mu\|_2\|S_\Lambda\|_1\gg N^{1/2}\cdot N^{1/2}=N$.
\end{enumerate}
\end{proposition}

\begin{proof}
(i) is Cauchy--Schwarz followed by
$\sum_{n\le N}\Lambda(n)^2\sim N\log N$ and
$\sum_{m\le N}\mu^2(m)\sim6N/\pi^2$. For (ii), the two factors are
bounded below for different reasons. Parseval gives
$\sup_\alpha|S_\mu|\ge\|S_\mu\|_2=(6N/\pi^2)^{1/2}$, an
identity-level constraint: no improvement in M\"obius
exponential-sum technology can lower $\sup_\alpha|S_\mu|$ below
$(6/\pi^2)^{1/2}N^{1/2}$. For the second factor Parseval runs the
\emph{wrong way}: it gives $\|S_\Lambda\|_1\le\|S_\Lambda\|_2
\ll(N\log N)^{1/2}$, an upper bound, while the trivial lower bound
$\|S_\Lambda\|_1\ge\bigl|\int_0^1S_\Lambda(\alpha)e(-n\alpha)
d\alpha\bigr|=\Lambda(n)$ reaches only $\log N$. The bound
$\|S_\Lambda\|_1\gg N^{1/2}$ that (ii) uses is a theorem of Vaughan
\cite{Vau88}.
\end{proof}

Davenport's uniform bound $S_\mu(\alpha)\ll_AN(\log N)^{-A}$ is
therefore useless here: it is a saving over $N$, whereas the pairing
needs a bound at the scale $N^{1/2}$, which Parseval forbids. This is
the parity obstruction in circle-method language --- \emph{the binary
problem sits exactly at the trivial bound with no margin} --- and it
is why the method that settles the ternary problem cannot be pushed to
the binary one by any sharpening of the M\"obius input.

\begin{measurement}[the constants]\label{meas:circle}
Computed by exact FFT on a $4N$-point grid at
$N=2^{14},\dots,2^{20}$:
\[
\begin{array}{r|cccc}
 N & 2^{14} & 2^{16} & 2^{18} & 2^{20}\\\hline
 \|S_\mu\|_2/\sqrt N & 0.7798 & 0.7797 & 0.7797 & 0.7797\\
 \sup|S_\mu|/\sqrt N & 3.058 & 2.742 & 2.853 & 2.801\\
 \|S_\Lambda\|_1/\sqrt N & 1.946 & 2.084 & 2.219 & 2.346\\
 \text{(i) bound}/N & 2.297 & 2.473 & 2.639 & 2.795\\
 N/\bigl(\sup|S_\mu|\,\|S_\Lambda\|_1\bigr) & 0.168 & 0.175 & 0.158
  & 0.152
\end{array}
\]
The first row reproduces $\sqrt{6/\pi^2}=0.7797$ exactly, confirming
the computation; the last row --- the margin that route~(ii) would
need to exceed $1$ --- sits near $0.16$ and falls over the last three
abscissae. Route~(i) diverges from the trivial bound like
$(\log N)^{1/2}$, as predicted. For reference the object itself is
small: $C(N)/N=-0.0105,\,0.0001,\,0.0059,\,0.0032$ at these $N$
(\texttt{audit\_circle\_margin.py}).
\end{measurement}

% =====================================================================
\section{Summary}

\begin{itemize}
\item The design space of divisor-switch weights is described exactly
by seven finite rearrangements
(Propositions~\ref{prop:dilate}--\ref{prop:direct}): the per-modulus
discrepancy is a dilate, the signs $\mu(k)$ cancel on the $\log k$
branch, and the untruncated object is the Goldbach count itself.

\item Theorem~\ref{thm:D}: no weight whose transform is
Bombieri--Vinogradov-accessible extracts $C(N)$; the loss is
$\exp(c_1\sqrt{(1/2)\log N})$, and the $1/2$ is the exponent of the
square-root barrier.

\item Theorem~\ref{thm:Dprime}: the separation survives
$EH(N^{\theta_E})$ for every fixed $\theta_E<1$. Closing it would
require $\theta_E=1$, where the statement is vacuous.

\item Proposition~\ref{prop:Dpp}: polynomial weights, which lie
outside the hypothesis of Theorem~\ref{thm:D}, are closed by
nonnegativity of $\Lambda_D$ and by circularity of the tuning.

\item Proposition~\ref{prop:E}: the circle method applied directly to
$C(N)$ has zero margin, with Parseval supplying the floor.

\item Net progress toward Goldbach: zero. What is closed is a design
space, so that it need not be re-explored.
\end{itemize}

% =====================================================================
\begin{thebibliography}{99}

\bibitem{Bom76} E.~Bombieri, \emph{The asymptotic sieve},
Rend. Accad. Naz. XL (5) \textbf{1/2} (1975/76), 243--269.

\bibitem{GY} D.~Goldston and C.~Y{\i}ld{\i}r{\i}m,
\emph{Higher correlations of divisor sums related to primes~I},
Integers \textbf{3} (2003), A5.

\bibitem{HL} J.-J.~Huang and H.~Li, \emph{On the connection between
the Goldbach conjecture and the Elliott--Halberstam conjecture},
arXiv:2005.03811v2 [math.NT], 2022.

\bibitem{Li23} J.~D.~Lichtman (with an appendix by S.~Drappeau),
\emph{Primes in arithmetic progressions to large moduli, and Goldbach
beyond the square-root barrier}, arXiv:2309.08522 [math.NT], 2023.

\bibitem{MV} H.~Montgomery and R.~Vaughan, \emph{Multiplicative number
theory~I: classical theory}, Cambridge University Press, 2007.

\bibitem{P1} \emph{An unconditional bound for a M\"obius-weighted
correlation sum in fixed residue classes, with two consequences for
the Huang--Li reduction of Goldbach to Elliott--Halberstam},
companion paper.

\bibitem{Vau88} R.~C.~Vaughan, \emph{The $L^1$ mean of exponential
sums over primes}, Bull. London Math. Soc. \textbf{20} (1988),
121--123.

\end{thebibliography}

% =====================================================================
\appendix
\section{Reproducibility}

\begin{center}
\begin{tabular}{lll}
\hline
Script & Result file & Supports\\
\hline
\texttt{lab\_dilate\_identity.py} & \texttt{lab\_dilate\_identity.txt}
 & Proposition~\ref{prop:dilate}\\
\texttt{lab\_positive\_weights.py} & \texttt{lab\_positive\_weights.txt}
 & Proposition~\ref{prop:posweights}\\
\texttt{lab\_weight\_gap.py} & \texttt{lab\_weight\_gap.txt}
 & Measurement~\ref{meas:flatsum}\\
\texttt{lab\_direct\_identity.py} & \texttt{lab\_direct\_identity.txt}
 & Proposition~\ref{prop:untrunc}\\
\texttt{lab\_layer\_decomposition.py} &
 \texttt{lab\_layer\_decomposition.txt} & Proposition~\ref{prop:layers}\\
\texttt{lab\_combined\_modulus.py} & \texttt{lab\_combined\_modulus.txt}
 & Proposition~\ref{prop:combined}\\
\texttt{lab\_direct\_route.py} & \texttt{lab\_direct\_route.txt}
 & Measurement~\ref{meas:direct}\\
\texttt{audit\_extraction\_tradeoff.py} &
 \texttt{audit\_extraction\_tradeoff.txt} & Measurement~\ref{meas:extract}\\
\texttt{audit\_polyweight.py} & \texttt{audit\_polyweight.txt}
 & Measurement~\ref{meas:poly}\\
\texttt{audit\_circle\_margin.py} & \texttt{audit\_circle\_margin.txt}
 & Measurement~\ref{meas:circle}\\
\hline
\end{tabular}
\end{center}

No computation is used as a premise of any proof. The identities of
Section~\ref{sec:identities} are finite rearrangements and are
verified as such to machine precision; the measurements record sizes
over the ranges stated in each statement.

\end{document}
